A finitely presented module is described by finitely many generators and relations.
These modules can be investigated by methods of linear algebra and play a central role in algebraic geometry.
When the base ring is coherent, this notion coincides with that of coherent modules, which is usually the generality one restricts to classically.

Synthetically, however, our base ring $R$ is provably not coherent~\cite[3.2]{topology-draft} and these two notions diverge.
As it turns out, neither is well-behaved in our setting.
The category of finitely presented $R$-modules lacks kernels, and the category of coherent $R$-modules, while always abelian, is much too small as it does not even contain crucial modules like $R^n$.
In this section, we study a larger category of finitely adequate modules~(\Cref{def:fin-adeq}) as a suitable replacement whose objects still retain a certain sense of finiteness.
This notion is not completely novel and has already been studied externally in~\cite[Chapter 06Z1]{stacks-project}.

The work in this section profited from discussions with Hugo Moeneclaey, Mohamed Barakat, and Marc Nieper-Wißkirchen.

\subsection{Finitely Adequate Modules Form an Abelian Category}

In this subsection, we will show that the category of finitely adequate $R$-modules, denoted by $\finadeq$, is abelian.
We begin by recounting the notion of quasi-coherence, introduced by Ingo Blechschmid~\cite{ingo-thesis}.
For this, recall that for an $R$-module $M$ and a finitely presented $R$-algebra $A$, there is an evaluation map:
\begin{align*}
  \ev : M \otimes A & \rightarrow M^{\Spec A} \\
  m \otimes a & \mapsto (f \mapsto f(a)m)
\end{align*}

\begin{definition}\label{def:quasi-coherent}
  Let $M$ be an $R$-module.
  $M$ is \notion{quasi-coherent (qc)} if for every finitely presented $R$-algebra $A$, the map $\ev$ is an equivalence.
  We say that $M$ is \notion{weakly quasi-coherent (wqc)} if this map is an equivalence for algebras of the form $R[r^{-1}]$, for each $r : R$.
\end{definition}

We can show that quasi-coherent $R$-modules $Q$ are tangential by a direct computation.

% \rednote{These pointed maps are already linear. Does linearity make sense when the domain is $\D(n)$? Respects scalar multiplication and addition whenever it exists.}

\begin{lemma}\label{lem:disk-qc}
  Let $Q$ be a quasi-coherent $R$-module.
  A pointed map $p : (Q,0)^{\D(n)}$ is an $n$-variate linear monomial with coefficients in $Q$.
\end{lemma}

\begin{proof}
  By quasi-coherence, we have the following equivalence $Q^{\D(n)} \simeq Q \otimes \quot{R[x_1,\dots,x_n]}{\ideal{x_ix_j}} \simeq Q^{n+1}$.
  We can thus identify the maps $\D(n) \rightarrow Q$ with $n$-variate polynomials with coefficients in $Q$ and which are of degree at most $1$.
  Let $p(x_1, \dots, x_n)$ be such a polynomial; then the canonical evaluation map acts like $\ev (p(x_1, \dots, x_n))(\varepsilon_1, \dots, \varepsilon_n) = p(\varepsilon_1, \dots, \varepsilon_n)$, where $\varepsilon_i : \D (n)$.
  Whenever $p$ represents a pointed map, we have $p(0,\dots,0) = p_0 = 0$.
\end{proof}

\begin{proposition}\label{prp:tangential-qc}
  Let $Q$ be a quasi-coherent $R$-module.
  Then $Q$ is tangential.
\end{proposition}

\begin{proof}
  By the above lemma, \Cref{lem:disk-qc}, we identify $\tangent{Q}{0}$ with univariate polynomials in $Q$ of degree $1$ and constant term $0$.
  Let $q : Q$ be a point in $Q$.
  Then $\can(q)(\varepsilon) = q\varepsilon$, where $\varepsilon : \D(1)$.
  We deduce that $\can(q) = qx$, so $\can$ is an equivalence by assumption.
\end{proof}

\begin{corollary}\label{cor:inf-lin-qc}
  Let $Q$ be a quasi-coherent $R$-module.
  Then $Q$ is infinitesimally linear at $0$.
\end{corollary}

\begin{proof}
  This follows immediately from \Cref{prp:tangential-qc} and \Cref{cor:tangential-inf-lin}.
\end{proof}

Since quasi-coherent modules are tangential, they form the basis of a well-behaved category of modules.
As tangentiality is preserved by finite limits and finite colimits, we will be especially interested in the kernel closure of this category.

\begin{definition}\label{def:adequate}
  An $R$-module $M$ is \notion{adequate} if it is the kernel of a linear map $f : P \rightarrow Q$ where $P$ and $Q$ are quasi-coherent $R$-modules.
\end{definition}

Adequate $R$-modules enjoy the following properties as they form the kernel closure of quasi-coherent modules.

\begin{proposition}\label{prp:wqc-adeq}
  Adequate $R$-modules are weakly quasi-coherent. \earlyqed
\end{proposition}

\begin{proposition}\label{prp:adeq-tangential}
  Adequate $R$-modules are tangential. \earlyqed
\end{proposition}

For the remainder of this section, we restrict our attention to a subcategory of adequate modules that satisfy certain finiteness conditions. We will construct the objects of this subcategory by first introducing their basic building blocks:

\begin{definition}\label{def:lin-adeq}
  Let $L$ be an $R$-module.
  $L$ is called \notion{linearly adequate} if it is the kernel of a linear map $f : R^m \rightarrow R^n$.
  We refer to the map $f$ as a \notion{linear copresentation} of $L$.
\end{definition}

\begin{remark}\label{rem:lin-adeq}
  In the article \cite{diffgeo-article}, linearly adequate modules went under the name finitely copresented modules.
  We decided to change the name here as linearly adequate matches the name of the external concept \cite[Chapter 06Z1]{stacks-project}.
\end{remark}

\begin{definition}\label{def:dual}
  Let $M$ be an $R$-module.
  We define the \notion{dual} of $M$ to be the module $M^\ast \defeq \Hom_R(M,R)$.
\end{definition}

\begin{theorem}[{{\cite[Lemma 2.3.9, Corollary 2.3.10]{diffgeo-article}}}]\label{thm:lin-adeq-dual}
  Dualization restricts to an anti-equivalence between linearly adequate modules and finitely presented ones.
  In other words, dualizing a linear copresentation of $L$,
  \[ M \hookrightarrow R^n \to R^m \]
  yields a finite presentation of $L^\ast$
  \[ R^m \to R^n \twoheadrightarrow L^\ast \]
  Moreover, the canonical map into the double dual, $\ev : L \to L^{\ast\ast}$, is an isomorphism for finitely presented and linearly adequate $R$-modules. \earlyqed
\end{theorem}

\begin{lemma}\label{lem:fp-lift}
  Let $M$ and $N$ be finitely presented $R$-modules and $f: M \to N$ a linear map.
  Then, there is an extension of this morphism to all presentations of $M$ and $N$:
  \begin{center}
    \begin{tikzcd}
      R^{m'}\ar[r]\ar[d,dotted,"\exists"] & R^m\ar[->>,r,"\pi_M"]\ar[d,dotted,"\exists"] & M\ar[d, "f"] \\
      R^{n'}\ar[r] & R^n\ar[->>,r,"\pi_N"] & N \\
    \end{tikzcd}
  \end{center}
  By duality, the analogous statement holds for linearly adequate modules.
\end{lemma}

\begin{proof}
  This is immediate since finite free $R$-modules are projective in $\modcat{R}$.
\end{proof}

From every commutative square, we can extract some exact sequences.
These will be particularly useful to show that finitely presented $R$-modules are closed under cokernels and linearly adequate $R$-modules are closed under kernels.

\begin{lemma}\label{lem:diff-ker-is-ker}
  Suppose we have a commutative square:
  \begin{center}
    \begin{tikzcd}
      A \ar{r}{u} \ar{d}{f} & B \ar{d}{g} \\
      C \ar{r}{v} & D
    \end{tikzcd}
  \end{center}
  The kernel $K$ of the induced map $\Ker(u) \rightarrow \Ker(v)$ is equivalent to the kernel $K'$ of $(u, f) : A \to B \oplus C$.
\end{lemma}

\begin{proof}
  Observe that $\Ker(u) \subseteq \Ker(gu)$.
  Thus $\Ker(u) = \Ker((u, gu) : A \to B \oplus D)$.
  Therefore, we obtain this result by the third isomorphism theorem:
  \begin{center}
    \begin{tikzcd}
      K \ar{r}{\simeq} \ar[hook]{d} & K' \ar[hook]{d} \\
      \Ker(u) \ar{d} \ar[hook]{r} & A \ar{d}{(u,f)} \ar{r}{(u,gu)} & B \oplus D \ar[equal]{d} \\
      \Ker(v) \ar[hook]{r} & B \oplus C \ar{r}{\mathrm{id} \oplus v} & B \oplus D
    \end{tikzcd}
  \end{center}
  \hfill{\phantom{.}}
\end{proof}

\begin{proposition}\label{prp:lin-adeq-ker-closure}
Linearly adequate $R$-modules are closed under kernels.
Dually, finitely presented $R$-modules are closed under cokernels.
\end{proposition}

\begin{proof}
  Suppose that $K$ is the kernel of a linear map $f : M \rightarrow N$, where $M$ and $N$ are linearly adequate.
  By the dual of \Cref{lem:fp-lift}, we can find a linear extension of $f$ to copresentations of $M$ and $N$:
  \begin{center}
    \begin{tikzcd}
      K \ar[hook]{d} \\
      M \ar{d} \ar[hook]{r} & R^{m_1} \ar{r} \ar{d} & R^{m_2} \ar{d} \\
      N \ar[hook]{r} & R^{n_1} \ar{r} & R^{n_2}
    \end{tikzcd}
  \end{center}
  By \Cref{lem:diff-ker-is-ker}, $K$ is also the kernel of a map $R^{m_1} \rightarrow R^{m_2 + n_1}$, so it is itself linearly adequate.
\end{proof}

\begin{lemma}\label{lem:ker-fp-is-cok-lin-adeq}
  If $f : M \to N$ is a map of finitely presented $R$-modules, then $\ker f$ is the cokernel of a map between linearly adequate $R$-modules.
\end{lemma}

\begin{proof}
  By \Cref{lem:fp-lift}, we can assume that $f$ is induced by a square of morphisms between finite free $R$-modules and construct the linearly adequate modules $A$ and $B$ as described below:
  \begin{center}
  \begin{tikzcd}
    A\ar[r,"\pi_1"]\ar[d,"l\times\id"] & R^{m'}\ar[r,"t"]\ar[d,"l"] & R^{n'}\ar[d,"r"]   \\
    B\ar[r,"\pi_1"]\ar[d] & R^m\ar[r,"b"]\ar[d] & R^n\ar[d] \\
    K\ar[r] &M\ar[r,"f"] & N
    \end{tikzcd}
   \end{center}
   $A:=\{(z,y):R^{m'}\times R^{n'}\mid blz=ry\}$, $B:=\{(x,y):R^{m}\times R^{n'}\mid bx=ry\}$.
\end{proof}

% We name the following modules after Adelman, since it will turn out that what we construct agrees with the free abelian category on the finite free $R$-modules, and our construction is close to Adelman's general construction of the free abelian category.

The modules in the kernel-cokernel closure of the finite free $R$-modules might yield a replacement of the algebro-geometric notion of coherent sheaf of modules, which would be a trivial notion over $R$.

\begin{definition}\label{def:fin-adeq}
  Let $M$ be an $R$-module. Then $M$ is a \notion{finitely adequate $R$-module} if one of the following equivalent statements holds:
  \begin{enumerate}[(i), noitemsep]
    \item $M$ is merely the kernel of a homomorphism between finitely presented $R$-modules.
    \item $M$ is merely the cokernel of a homomorphism between linearly adequate $R$-modules.
  \end{enumerate}
  We denote the type of finitely adequate modules with $\finadeq$\index{$\finadeq$}, as well as the category of these modules.
\end{definition}

\begin{proposition}\label{prp:fin-adeq-adeq}
  Every finitely adequate $R$-module is adequate.
  In particular, finitely adequate $R$-modules are also tangential and wqc.
\end{proposition}

\begin{proof}
  Every finitely presented $R$-module is quasi-coherent.
  Since every finitely adequate $R$-module is merely the kernel of a map between quasi-coherent modules, it is itself adequate.
  The rest follows from \Cref{prp:tangential-qc} and \Cref{prp:wqc-adeq}.
\end{proof}

We will first characterize the projective and injective objects in $\finadeq$.

\begin{lemma}\label{lem:lin-adeq-lift}
  Let $L$ be a linearly adequate $R$-module.
  Every lifting problem
  \begin{center}
    \begin{tikzcd}
      & A \ar[two heads]{d}{e} \\
      L \ar{r}{v} & B
    \end{tikzcd}
  \end{center}
  merely has a linear solution when $A$ and $B$ are tangential and $M \defeq \Ker e$ is wqc.
\end{lemma}

\begin{proof}
  By \cite[Theorem 7.3.6]{draft} we have $H^1(L,M)=0$, since $L$, like any linearly adequate $R$-module, is an affine scheme. Indeed, as the kernel of a linear map $f : R^m \to R^n$, the module $L$ is equivalent to $\Spec(\quot{R[x_1, \dots, x_m]}{\langle f(x_1), \dots, f(x_n) \rangle})$.

  By definition, we have $H^1(L,M)=||L \to K(M,1)||_0$ and therefore $L \to K(M,1)$ is connected.
  For any $b : B$, we let $E_b$ denote the fiber of $e$ over $b$.
  Then $E_b$ is an $M$-torsor, or in other words an element of $K(M,1)$.
  So $E \circ v$ defines a map $L \to K(M,1)$, and we have $||E \circ v = (\argmhole \mapsto M)||$ by the connectedness noted above.
  This means that the dependent type $E \circ v$ merely has a section, since $(\argmhole \mapsto M)$ has a section.
  Let $h'' : L \rightarrow A$ denote the section after projecting onto $A$.
  The map $h' \defeq h'' - h''(0)$ is now a pointed map of type $(L,0) \rightarrow_\ast (A,0)$.
  The map $h'$ still solves the problem as $e(h''(0)) = 0$.
  By \Cref{thm:linearize-lift}, there is a linear map $h : L \rightarrow A$ solving this lifting problem.
\end{proof}

\begin{theorem}\label{thm:lin-adeq-proj}
  Let $M$ be a finitely adequate module.
  Then $M$ is projective in $\finadeq$ if and only if $M$ is linearly adequate.
\end{theorem}

\begin{proof}
  Assume first that $M$ is linearly adequate.
  We need to merely solve any lifting problem of the form
  \begin{center}
    \begin{tikzcd}
      & A \ar[two heads]{d}{e} \\
      M \ar{r}{v} & B
    \end{tikzcd}
  \end{center}
  where $A$ and $B$ are finitely adequate.
  Such a solution follows from \Cref{lem:lin-adeq-lift}, since $A$ and $B$ are tangential and wqc by \Cref{prp:fin-adeq-adeq}.

  Assume now that $M$ is projective.
  Let $L_1 \xrightarrow{p} L_2 \twoheadrightarrow M$ be a presentation of $M$ by linearly adequate modules, and let $M' \defeq \Im p$.
  Since $M$ is assumed projective, the epimorphism $L_2 \twoheadrightarrow M$ splits.
  It follows that the inclusion $M' \hookrightarrow L_2$ splits as well and $M'$ is a direct summand of $L_2$.
  Since $L_2$ is linearly adequate, it is projective, thus $M'$ is projective as well.
  Since $M'$ is projective, the canonical map $L_1 \twoheadrightarrow M'$ splits.
  This gives us a map $L_2 \rightarrow L_1$ such that $M$ is its kernel.
  $M$ is now linearly adequate by \Cref{prp:lin-adeq-ker-closure}, as it is the kernel of a map between linearly adequate $R$-modules.
\end{proof}

\begin{lemma}\label{lem:fin-adeq-lift}
  Any linear map $f : M \rightarrow N$ merely lifts to a commutative square of linearly adequate modules.
\end{lemma}

\begin{proof}
  This is a direct consequence of \Cref{thm:lin-adeq-proj}.
\end{proof}

\begin{proposition}\label{prp:fin-adeq-cok-closure}
  Finitely adequate modules are closed under cokernels.
\end{proposition}

\begin{proof}
  This proof follows from \Cref{lem:diff-ker-is-ker} and \Cref{lem:fin-adeq-lift}, following a proof analogous to \Cref{prp:lin-adeq-ker-closure}.
\end{proof}

\begin{lemma}\label{lem:dual-exact-fp-inc}
  Let $M$ be a finitely adequate $R$-module, and $P$ a finitely presented $R$-module.
  Any lifting problem
  \begin{center}
    \begin{tikzcd}
      M \ar[hook]{d}{m} \ar{r}{u} & R \\
      P
    \end{tikzcd}
  \end{center}
  merely admits a linear solution.
\end{lemma}

\begin{proof}
  By \Cref{prp:fin-adeq-cok-closure}, the cokernel $N \defeq \Cok(m)$ is finitely adequate.
  Thus, there exists a monomorphism $n : N \hookrightarrow Q$ such that $Q$ is finitely presented.
  Let $f : P \rightarrow Q$ be the map induced by $n$.
  Now, $m$ is the kernel of $f$.
  By \Cref{lem:ker-fp-is-cok-lin-adeq}, we obtain an analysis of $M$:
  \begin{center}
    \begin{tikzcd}
      K \ar{d}{p \times \id} \ar[hook]{r} & R^{m_1} \ar{r}{f_1} \ar{d}{p} & R^{n_1} \ar{d}{q} \\
      L \ar[two heads]{d} \ar[hook]{r} & R^{m_2} \ar{r}{f_2} \ar[two heads]{d} & R^{n_2} \ar[two heads]{d} \\
      M \ar[hook]{r}{m} & P \ar{r}{f} & Q
    \end{tikzcd}
  \end{center}
  In the diagram above, we have $P = \Cok(p)$, $Q = \Cok(q)$, $K = R^{m_1} \underset{R^{n_2}}{\times} R^{n_1}$, and $L = R^{m_2} \underset{R^{n_2}}{\times} R^{n_1}$.
  Let $t$ be the composite $f_2 \circ p = q \circ f_1$.
  
  In the above analysis, only the columns are exact.
  We will alter this diagram so that every column and row becomes exact.
  \begin{center}
    \begin{tikzcd}
      K \ar{d}{p \times \id} \ar[hook]{r} & R^{m_1 + n_1} \ar{r}{t + q} \ar{d}{p \times \id} & R^{n_2} \ar[equal]{d} \\
      L \ar[two heads]{d} \ar[hook]{r} & R^{m_2 + n_1} \ar{r}{f_2 + q} \ar[two heads]{d} & R^{n_2} \\
      M \ar[hook]{r}{m} & P
    \end{tikzcd}
  \end{center}

  The linear map $u : M \to R$ restricts to a map $l_1 : L \to R$ that is zero on $K$.
  By~\cite[Lemma 2.3.7]{diffgeo-article}, it is possible to extend $l_1$ to $l_2 : R^{m_2+n_1} \to R$.
  $l_2$ will descend to $M$ if and only if it vanishes on the image of $p \times \id$.
  Although this is not generally the case, we can construct a linear correction $c : R^{m_2} \to R$ such that $l_2 - c$ has this property.

  To construct $c$, let $c_0 \coloneq l_2 \circ (p \times \id)$.
  Note that $c_0$ vanishes on $K$ and therefore descends to a linear map $c_1 : \Im(t + q) \to R$.
  By \cite[Lemma 2.3.8]{diffgeo-article}, $c_1$ extends to a map $c_2 : R^{n_2} \to R$.
  Now, $c \coloneq c_2 \circ (f_2 + q)$ is a linear map on $R^{m_2 + n_1}$ that is zero on $K$,
  so $l_2 - c$ vanishes on the image of $p \times \id$.
  Therefore, $l_2 - c$ descends to a map $l : P \to R$.

  Finally, note that $l$ restricts to $l_1$ on $L$.
  Thus we have $u = l \circ m$.
\end{proof}

Let $M$ be an $R$-module.
There is a natural map $\ev : M \rightarrow M^{\ast\ast}$ that sends a term $m$ to $\ev_m(f) \defeq f(m)$.

\begin{theorem}\label{thm:can-dual-iso}
  Let $M$ be a finitely adequate module.
  Then $\ev : M \rightarrow M^{\ast\ast}$ is an equivalence.
\end{theorem}

\begin{proof}
  We will prove this by exploiting two different naturality squares involving $\ev$.

  Since $M$ is finitely adequate, it merely embeds into a finitely presented $R$-module $P$.
  This yields the naturality square
  \begin{center}
    \begin{tikzcd}
      M \ar[hook]{r} \ar{d}[right]{\ev} \ar[hook]{rd} & P \ar{d}[right]{\ev}[left]{\simeq} \\
      M^{\ast\ast} \ar{r} & P^{\ast\ast}
    \end{tikzcd}
  \end{center}
  Since the composite $M \rightarrow P^{\ast\ast}$ is a monomorphism, so $\ev : M \rightarrow M^{\ast\ast}$ must be a monomorphism as well.

  Since $M$ is finitely adequate, it is merely covered by a linearly adequate $R$-module $L$.
  Observe that the dual of the cover, $M^\ast \hookrightarrow L^\ast$, is a monomorphism.
  Taking the dual again yields an epimorphism $L^{\ast\ast} \twoheadrightarrow M^{\ast\ast}$, since $L^{\ast}$ is finitely presented by \Cref{thm:lin-adeq-dual} and \Cref{lem:dual-exact-fp-inc}.
  This yields the naturality square
  \begin{center}
    \begin{tikzcd}
      L \ar[two heads]{r} \ar{d}[right]{\ev}[left]{\simeq} \ar[two heads]{rd} & M \ar{d}{\ev} \\
      L^{\ast\ast} \ar[two heads]{r} & M^{\ast\ast}
    \end{tikzcd}
  \end{center}
  Since the composite $L \rightarrow M^{\ast\ast}$ is an epimorphism, $\ev : M \rightarrow M^{\ast\ast}$ has to be an epimorphism as well.
\end{proof}

\begin{corollary}\label{cor:dual-equiv}
  Taking duals $(\argmhole)^{\ast} : \finadeq \rightarrow \finadeq$ is an anti-equivalence of categories. \earlyqed
\end{corollary}

We get the following results following duality.

\begin{theorem}\label{thm:fp-inj}
  Let $M$ be finitely adequate.
  Then $M$ is injective in $\finadeq$ if and only if $M$ is finitely presented.
\end{theorem}

\begin{proof}
  By \Cref{cor:dual-equiv} it is sufficient to prove the dual statement: $M$ is projective in $\finadeq$ if and only if $M$ is linearly adequate.
  This is precisely \Cref{thm:lin-adeq-proj}.
\end{proof}

\begin{corollary}\label{cor:fin-adeq-extension}
  Any linear map $f : M \rightarrow N$ merely extends to a commutative square of finitely presented $R$-modules. \earlyqed
\end{corollary}

\begin{corollary}\label{cor:fin-adeq-ker-closure}
  Finitely adequate modules are closed under kernels. \earlyqed
\end{corollary}

Since the category of finitely adequate modules is a subcategory of $\modcat{R}$ which is closed under finite limits and finite colimits, it is itself abelian.

\begin{theorem}\label{thm:fin-adeq-abelian}
  The category $\finadeq$ is abelian. \earlyqed
\end{theorem}

\begin{theorem}\label{thm:fin-adeq-class}
  Let $M$ be an $R$-module.
  The following are equivalent:
  \begin{enumerate}[(i), noitemsep]
    \item $M$ is finitely adequate.
    \item $M$ is the image of a map from a linearly adequate $R$-module to a finitely presented $R$-module.
    \item $M$ is in the smallest abelian subcategory of $\modcat{R}$ that contains all finite free $R$-modules.
  \end{enumerate}
\end{theorem}

\begin{proof}
  The equivalence between (i) and (ii) is clear from \Cref{prp:fin-adeq-cok-closure} and \Cref{cor:fin-adeq-ker-closure}.
  The equivalence with (iii) will follow from \Cref{thm:fin-adeq-abelian}.
\end{proof}

\begin{remark}\label{rem:fin-adeq-class-self-dual}
  Observe that condition (ii) is self-dual.
\end{remark}

\begin{remark}\label{rem:fin-adeq-Adelman}
  By \cite[Theorem 1.15]{adelman-construction}, \Cref{thm:lin-adeq-proj} and \Cref{thm:fp-inj} imply that the category $\finadeq$ is the free abelian category over the finite free $R$-modules.
  This means the finitely adequate $R$-modules have the following universal property:
  For every abelian category $\mathcal{A}$ and all additive functors $F:\finfree \to \mathcal{A}$, there exists a unique exact functor $K$ as follows:
  \begin{center}
    \begin{tikzcd}
      \finfree \ar{d} \ar{r}{F} & \mathcal{A} \\
      \finadeq \ar[dashed]{ru}[swap]{K}
    \end{tikzcd}
  \end{center}
  Consequently, the data needed to specify an exact functor $K$ is an assignment $KR$ of $R$ in $\mathcal{A}$ and a group homomorphism $\Hom_R(R,R) \rightarrow \mathcal{A}(KR,KR)$.
\end{remark}

\subsection{Properties of Finitely Adequate Modules}

In this subsection we collect some properties of finitely adequate $R$-modules.

\begin{proposition}
  Let $A$ be a closed submodule of a finitely adequate module $M$. Then $A$ is finitely adequate if $A$ is tangential.
\end{proposition}

\begin{proof}
  Let $M$ be presented by the linearly adequate modules $K \to L \twoheadrightarrow M$. Consider the pullback $L' \coloneq A \times_M L$.
  Since pullbacks of $R$-modules are taken as pullbacks on the underlying sets, we have that $L'$ is a closed submodule of $L$, and is therefore a scheme by \cite[Proposition 5.4.3]{draft}.
  Therefore the tangent space $\tangent{L'}{0}$ is a linearly adequate module by \cite[Lemma 2.2.8]{diffgeo-article}.
  Since tangential modules are closed under limits \Cref{prp:op-closure-inf-lin-tangential}, $L'$ is also tangential, and therefore also linearly adequate.
  The sequence $K \to L' \twoheadrightarrow A$ presents $A$ as a cokernel of a map of linearly adequate modules, so $A$ is finitely adequate.
\end{proof}

\begin{proposition}\label{prp:fin-adeq-hom-closure}
  For finitely adequate $R$-modules $M$ and $N$, $\Hom_R(M,N)$ is a finitely adequate $R$-module as well.
\end{proposition}

\begin{proof}
  The proof in \cite[Chapter IV, 4.12]{lombardi-quitte} works, if ``coherent'' is replaced by ``finitely adequate''.

  First we use their proof to show that for finitely presented $R$-modules $M$ and $N$, the $R$-module $\Hom_R(M,N)$ is finitely adequate. Morphisms $f: M \to N$ are presented by squares:
  \begin{center}
    \begin{tikzcd}
      R^m \ar{r} \ar{d}{\varphi} & R^{m'} \ar{d}{\psi} \\
      R^n\ar{r} & R^{n'}
    \end{tikzcd}
  \end{center}
  So we have a finite free module of pairs of morphisms $(\varphi,\psi)$ and a submodule of pairs such that the square above commutes.
  This submodule $S$ is the kernel of a linear map and therefore finitely adequate.
  We have a surjection $\pi : S \to \Hom_R(M,N)$, so $\Hom_R(M,N)$ is finitely adequate if $\pi$ is a cokernel of a map of finitely adequate $R$-modules.
  This is the case, since there is a surjection onto $\Ker{\pi}$ from the finite free $R$-module of linear maps which split the square:
  \begin{center}
    \begin{tikzcd}
      R^m \ar{r} \ar{d}{\varphi} & R^{m'} \ar{d}{\psi} \ar{ld}{s} \\
      R^n \ar{r} & R^{n'}
    \end{tikzcd}
  \end{center}
  Then we can reuse the same argument to show the statement of the proposition.
\end{proof}

\begin{proposition}\label{prop:hood-faithful}
  The functor $\hood{\argmhole}{1} : \modcat R \rightarrow \pType$ is faithful when restricted to the subcategory of tangential $R$-modules.
\end{proposition}

\begin{proof}
  Let $f,g : M \rightrightarrows N$ be two linear maps between tangential $R$-modules with $\hood{f,0}{1} = \hood{g,0}{1}$.
  Then $df_0=dg_0$ and by naturality of the map $\can$ the claim follows.
  % This is immediate, as the action on the map is given by restriction, i.e. $\hood{f,0}{1} \defeq f\mid_{\hood{\argmhole}{1}}$.
  % To see that it is full, let $M$ and $N$ be finitely adequate $R$-modules, and let $f : \hood{M,0}{1} \rightarrow_\ast \hood{N,0}{1}$ be a pointed map.
  % Because $M$ and $N$ are finitely adequate, we can find linearly adequate $R$-modules $K_1$, $K_2$, $L_1$, and $L_2$ together with presentations $K_1 \rightarrow K_2 \twoheadrightarrow M$ and $L_1 \rightarrow L_2 \twoheadrightarrow N$.
  % By ???, $\hood{K_i,0}{1}$ and $\hood{L_i,0}{1}$ are projective types.
  % Using this property, we can lift the map $f$ to a commutative square by \Cref{lem:hood-module-prs-surj}
  % \begin{center}
  %   \begin{tikzcd}
  %     \hood{K_1,0}{1} \ar{d} \ar[dashed]{r} & \hood{L_1,0}{1} \ar{d} \\
  %     \hood{K_2,0}{1} \ar[two heads]{d} \ar[dashed]{r} & \hood{L_2,0}{1} \ar[two heads]{d} \\
  %     \hood{M,0}{1} \ar{r}{f} & \hood{N,0}{1}
  %   \end{tikzcd}
  % \end{center}
  % By \cite[Lemma 2.1.4]{diffgeo-article} and \Cref{thm:lin-adeq-dual}, the morphisms in the top square extend to modules.
  % By taking cokernels we get an induced map
  % \begin{center}
  %   \begin{tikzcd}
  %     K_1 \ar{r} \ar{d} & L_1 \ar{d} \\
  %     K_2 \ar[two heads]{d} \ar{r} & L_2 \ar[two heads]{d} \\
  %     M \ar{r}{f'} & N
  %   \end{tikzcd}
  % \end{center}
  % It is clear that $f'\mid_{\hood{M,0}{1}} = f$.
\end{proof}

\begin{remark}
  The functor $\hood{\argmhole}{1}$ of the previous proposition is full when restricted even further to linearly adequate $R$-modules \rednote{cite}.
  This, however, already fails on finitely adequate $R$-modules.

  To see this, consider the maps $\alpha_\varepsilon : (\varepsilon) \to (\varepsilon), r\varepsilon \mapsto r^2\varepsilon$ for every $\varepsilon:\D(1)$.
  These are well-defined, since $r\varepsilon = s\varepsilon$ implies $r^2\varepsilon-s^2\varepsilon = (r+s)(r-s)\varepsilon = 0$.
  Since for every $n:(\varepsilon)$ we have $(\varepsilon=0)\to(n=0)$, we conclude that $(\varepsilon) = \hood{(\varepsilon),0}{1}$.
  Note that $\hood{\argmhole}{1}$ is the identity functor on those $R$-modules that are their own first order neighbourhood.
  Thus it suffices to show that it is not the case that all $\alpha_\varepsilon$ are $R$-linear.
  Assume they are.
  Then $r^2\varepsilon = \alpha_\varepsilon(r\varepsilon) = r\alpha_\varepsilon(\varepsilon) = r\varepsilon$ and hence also $(r^2-r)\varepsilon=0$ for all $r:R$.
  Using duality and comparing coefficients, we see that $\varepsilon=0$ for every $\varepsilon:\D(1)$, which cannot be.
\end{remark}

\rednote{The next remark could maybe move to the appendix.  Not sure if we need it in the final paper.}

\begin{remark}\label{rem:hood-module-prs-surj}
  If $f : M \rightarrow N$ be an epimorphism of $R$-modules with $M$ linearly adequate
  and $N$ finitely adequate,
  it does not follow that $\hood{f,0}{1}$ is surjective.
  Here is an example.

  Let $\varepsilon : \D(1)$ be arbitrary.
  Consider the map $\pi_\varepsilon : R \to (\varepsilon)$ which sends $r$ to $r \varepsilon$.
  Note that $R$ is linearly adequate and $(\varepsilon)$ is finitely adequate
  by \Cref{thm:fin-adeq-class}.
  We will show that $\pi_\varepsilon$ does not satisfy the claim for all $\varepsilon$.
  Indeed, assume that $\hood{\pi_{\varepsilon}, 0}{1}$ is surjective for all $\varepsilon$.
  Note that $\varepsilon \in \hood{(\varepsilon), 0}{1} = (\varepsilon)$.
  Let $r:\D(1) = \hood{R,0}{1}$ be a lift of $\varepsilon$.
  Since $\pi_\varepsilon(1) = \varepsilon$, $\pi_\varepsilon(r-1)=0$, i.e., $(r-1)\varepsilon=0$.
  But $r$ is nilpotent, hence $r-1$ invertible and this would mean $\varepsilon = 0$, which cannot be true for all $\varepsilon$.
\end{remark}

Finitely adequate modules form a weak Serre subcategory of $\modcat{R}$.

\begin{lemma}\label{lem:fin-adeq-ext-pb-red}
  Suppose we have a commutative diagram
  \begin{center}
    \begin{tikzcd}
      B \ar[equal]{d} \ar[hook]{r} & p^*E \ar[two heads]{r} \ar[two heads]{d} \ar[phantom]{rd}[very near start]{\lrcorner}[very near end]{\ulcorner} & P \ar[two heads]{d}{p} \\
      B \ar[hook]{r} & E \ar[two heads]{r} & A
    \end{tikzcd}
  \end{center}
  in which $B$, $A$ and $P$ are finitely adequate $R$-modules.
  Then $E$ is finitely adequate if and only if $p^*E$ is so.
\end{lemma}

\begin{proof}
  Suppose that $E$ is finitely adequate.
  Since $p^*E$ is a pullback of finitely adequate $R$-modules, it is itself finitely adequate.

  Now assume that $p^*E$ is finitely adequate.
  By the third isomorphism theorem we can prove that $E$ is finitely adequate.
  In the following diagram, the object $0^*_Ap^*E$ is the kernel of $p^*E \rightarrow E$, $0^*_AE$ is the kernel of $E \rightarrow A$, and $0^*_Ep^*E$ is the kernel of $p^*E \rightarrow E$.
  \begin{center}
    \begin{tikzcd}
      0^*_Ep^*E \ar[equal]{r} \ar[hook]{d} & 0^*_Ep^*E \ar[hook]{d} \\
      0^*_Ap^*E \ar[hook]{r} \ar[two heads]{d} & p^*E \ar[two heads]{r} \ar[two heads]{d} & A \ar[equal]{d} \\
      B \ar[hook]{r} & E \ar[two heads]{r} & A
    \end{tikzcd}
  \end{center}
  Chasing around this diagram we can see that $0^*_Ap^*E$ is finitely adequate, because $p^*E$ is.
  Furthermore, $0^*_Ep^*E$ is finitely adequate, because $0^*_Ap^*E$ is finitely adequate.
  It then follows that $E$ is finitely adequate, because it is the cokernel of two finitely adequate $R$-modules.
\end{proof}

\begin{lemma}\label{lem:fin-adeq-ext-red}
  Finitely adequate $R$-modules are closed under extension if and only if for any linearly adequate $R$-module $L$ and finitely adequate $R$-module $B$, any exact sequence $0 \to B \to E \to L \to 0$ splits.
\end{lemma}

\begin{proof}
  Assume that finitely adequate modules are closed under extensions.
  Then it is clear that any extension by $B$ and $L$ is split, as $L$ is injective among the finitely adequate modules \Cref{thm:lin-adeq-proj}.

  Assume now that $\textup{Ext}^1(L,B) = 0$ for any such modules $L$ and $B$.
  Consider an extension problem
  \begin{center}
    \begin{tikzcd}
      B \ar[hook]{r} & E \ar[two heads]{r} & A
    \end{tikzcd}
  \end{center}
  where $B$ is finitely adequate, and $L$ is finitely adequate and is merely presented by $L_1 \rightarrow L_2 \overset{p}{\twoheadrightarrow} A$.
  By \Cref{lem:fin-adeq-ext-pb-red}, we can take the pullback of the extension along $p : P_1 \twoheadrightarrow L$ to get another associated extension.
  \begin{center}
    \begin{tikzcd}
      B \ar[equal]{d} \ar[hook]{r} & p^*E \ar[two heads]{d} \ar[two heads]{r} \ar[phantom]{rd}[very near start]{\lrcorner} & L_1 \ar[two heads]{d} \\
      B \ar[hook]{r} & E \ar[two heads]{r} & A
    \end{tikzcd}
  \end{center}
  Now $E$ is finitely adequate if and only if $p^*E$ is so.

  Since we assumed that any exact sequence on the form $0 \to B \to E \to L_1 \to 0$ splits, it is necessary that $p^*B = B \oplus L_1$, which is finitely adequate.
  Thus, $E$ is finitely adequate as desired.
\end{proof}

\begin{theorem}\label{thm:fin-adeq-ext-closed}
  Finitely adequate $R$-modules are closed under extension.
\end{theorem}

\begin{proof}
  By \Cref{lem:fin-adeq-ext-red} it is sufficient to show that any exact sequence on the form $0 \to B \to E \to L \to 0$ splits for any finitely adequate $B$ and linearly adequate $L$.
  This amounts to showing that any lifting problem
  \begin{center}
    \begin{tikzcd}
      & E \ar[two heads]{d}{p} \\
      L \ar[equal]{r} & L
    \end{tikzcd}
  \end{center}
  where $L$ is finitely adequate and $B \defeq \Ker(p)$ is finitely adequate, merely has a linear solution.
  Observe that $E$ is tangential by \Cref{prp:op-closure-inf-lin-tangential}, since $B$ and $L$ are tangential by \Cref{prp:adeq-tangential}.
  $B$ is also wqc by \Cref{prp:wqc-adeq}.
  Thus the lifting problem has a solution by \Cref{lem:lin-adeq-lift}.
\end{proof}


By \Cref{rem:tensor-not-closed}, the tensor product $M \otimes N$ is not necessarily linearly adequate when $M$ and $N$ are linearly adequate.
This contrasts with the case where $M$ is finitely presented (or, similarly, where $N$ is finitely presented), since $M \otimes N$ can then be realized as the cokernel of a map between finite direct sums of $N$, and is therefore finitely presented.
Somewhat surprisingly, the internal hom still has a left adjoint in $\finadeq$.
We will show that the \notion{double dual tensor product} $M \dbldualtensor N \defeq (M \otimes N)^{\ast\ast}$ is this left adjoint.
By the hom-tensor adjunction, there is already a natural isomorphism $\Hom_R(A, \Hom_R(B, C)) \simeq \Hom_R(A \otimes B, C)$.
Thus, it is sufficient to prove that $\ev_{A \otimes B}^\ast : \Hom_R(A \dbldualtensor B, C) \rightarrow \Hom_R(A \otimes B, C)$ is a natural isomorphism for all finitely adequate modules $A$, $B$, and $C$.
The composite of these natural isomorphisms yields the natural isomorphism $\Hom_R(A \dbldualtensor B, C) \simeq \Hom_R(A, \Hom_R(B, C))$.
It is clear that $\ev_{A \otimes B}^\ast$ is natural in all three arguments.
Now, $\ev_{A \otimes B}^\ast$ is an isomorphism precisely when every morphism $f : A \otimes B \to C$ factors uniquely through $\ev_{A \otimes B} : A \otimes B \to A \dbldualtensor B$.

We first carry out this factorization in the general setting of monads.
Let $T$ be a monad on $\mathcal{C}$ with multiplication $\mu$ and unit $\eta$.
In our situation, $\mathcal{C}$ will be the category of $R$-modules and $T$
will be the double-dual operation.

\begin{definition}\label{def:monad-reflexive}
  Let $X$ be an object of $\mathcal{C}$. We say that $X$ is:
  \begin{itemize}
    \item \notion{$T$-torsionless} if the unit $\eta_X$ is a monomorphism,
    \item \notion{$T$-reflexive} if the unit $\eta_X$ is an isomorphism.
  \end{itemize}
\end{definition}

\begin{lemma}\label{lem:monad-app-torsionless}
  Let $X$ be an object of $\mathcal{C}$. Then $TX$ is $T$-torsionless.
\end{lemma}

\begin{proof}
  By unitality, we have $\mu_X \circ \eta_{TX} = \id_{TX}$, so $\eta_{TX}$ is a split monomorphism.
\end{proof}

\begin{definition}\label{def:rel-epi}
  Let $\mathfrak{C}$ be a class of objects in $\mathcal{C}$. A map $f : X \to Y$ is \textbf{$\mathfrak{C}$-epic} if for every $Z$ in $\mathfrak{C}$ and all morphisms $g, h : Y \to Z$ such that $g \circ f = h \circ f$, we have $g = h$.
\end{definition}

The monad $T$ is said to be \notion{well-pointed} if $T\eta_X = \eta_{TX}$ for every object $X$ of $\mathcal{C}$.
We will not discuss well-pointed monads in this global sense, but rather focus on the local notion.
For a treatment of well-pointed monads, we refer the reader to Chapter 4.2 of Borceux's book~\cite{Borceux_1994_II}.
The next result is similar to Proposition 4.2.3 of Borceux~\cite{Borceux_1994_II}, but we do not assume that $T$ is well-pointed, which forces us to give a different proof.

\begin{proposition}\label{prp:monad-well-pointedness-torsionless-epi}
  Let $X$ be an object of $\mathcal{C}$. The following are equivalent:
  \begin{enumerate}[(i), noitemsep]
    \item $TX$ is $T$-reflexive,
    \item $T\eta_X = \eta_{TX}$,
    \item $\eta_X$ is $T$-torsionless-epic.
  \end{enumerate}
\end{proposition}

\begin{proof}
  (i) $\Rightarrow$ (iii): Assume that $TX$ is $T$-reflexive, i.e., that $\eta_{TX}$ is an isomorphism.
  Then $T\eta_X$ is an isomorphism as well, since $T\eta_X$ is a section of $\mu_X$.
  Consider two morphisms $f, g : TX \to Y$, where $Y$ is $T$-torsionless and $f \circ \eta_X = g \circ \eta_X$.
  By applying $T$ to this equality, we obtain $Tf \circ T\eta_X = Tg \circ T\eta_X$.
  Thus, $Tf = Tg$ since $T\eta_X$ is an epimorphism.
  By naturality of $\eta$, we have $\eta_Y \circ f = Tf \circ \eta_{TX} = Tg \circ \eta_{TX} = \eta_Y \circ g$.
  Since $Y$ is $T$-torsionless, $\eta_Y$ is a monomorphism and thus $f = g$.
  
  (iii) $\Rightarrow$ (ii):  Assume that $\eta_X$ is $T$-torsionless-epic.
  By naturality of $\eta$, we have $\eta_{TX} \circ \eta_X = T\eta_X \circ \eta_X$. 
  Since $TX$ is $T$-torsionless by \Cref{lem:monad-app-torsionless}, we obtain $T\eta_X = \eta_{TX}$.
  
  (ii) $\Rightarrow$ (i): Assume that $T\eta_X = \eta_{TX}$.
  Our goal is to show that $\eta_{TX}$ is an isomorphism, so it suffices to show that $T\eta_X$ is an isomorphism.
  By unitality, we have that $\mu_X \circ T\eta_X = \id_{TX}$, so it remains to check the other composite.
  By unitality and using $T\eta_X = \eta_{TX}$, we have
  \[
    \id_{T^2X} = \mu_{TX} \circ T\eta_{TX} = \mu_{TX} \circ T^2\eta_{X} .
  \]
  From the naturality square of $\mu$ applied to $\eta_X$, we get the following commutative square:
  \begin{center}
    \begin{tikzcd}
      T^2X \ar{d}{\mu_X} \ar{r}{T^2\eta_{X}} & T^3X \ar{d}{\mu_{TX}} \\
      TX \ar{r}{T\eta_{X}} & T^2X
    \end{tikzcd}
  \end{center}
  Combining these two facts, we deduce that $T\eta_X \circ \mu_X = \id_{T^2X}$, so $T\eta_X$ is an isomorphism.
\end{proof}


\begin{proposition}\label{prp:monad-factorization}
  Let $f : X \to Y$ be a morphism in $\mathcal{C}$. Then $f$ factors uniquely through $\eta_X$ if both $TX$ and $Y$ are $T$-reflexive.
\end{proposition}

\begin{proof}
  The naturality square of $\eta$ applied to $f$
\begin{center}
  \begin{tikzcd}
    X \ar{r}{f} \ar{d}{\eta_X} & Y \ar{d}{\eta_Y}[left]{\simeq} \\
    T X \ar{r}{T f} & T Y
  \end{tikzcd}
\end{center}
  yields a factorization $f = \eta_Y^{-1} \circ T f \circ \eta_X$ of $f$ through $\eta_X$.
  For uniqueness, assume we are given $f', f'' : TX \to Y$ such that $f' \circ \eta_X = f'' \circ \eta_X$.
  By \Cref{prp:monad-well-pointedness-torsionless-epi}, $\eta_X$ is $T$-torsionless-epic, which yields $f' = f''$.
\end{proof}

Now we specialize to a monad coming from an adjoint pair
$F \dashv G$ with $F : \mathcal{C} \to \mathcal{D}$, $G : \mathcal{D} \to \mathcal{C}$, unit $\eta : \id_{\mathcal{C}} \to GF$ and counit $\varepsilon : FG \to \id_{\mathcal{D}}$.
The associated monad is $T \coloneq GF$, with multiplication $\mu \coloneq G \varepsilon_F$ and unit $\eta$.
In the case where the monad is induced from an adjunction, we can make a slight improvement to \Cref{prp:monad-well-pointedness-torsionless-epi}.

\begin{proposition}\label{prp:monad-adjunction-well-pointed-improvement}
  Let $X$ be an object of $\mathcal{C}$.
  Then $GF\eta_X = \eta_{GFX}$ if and only if $F\eta_X$ is an isomorphism.
\end{proposition}

\begin{proof}
  Assume first that $GF\eta_X = \eta_{GFX}$.
  This part of the proof is completely analogous to the implication (ii) $\Rightarrow$ (i) in \Cref{prp:monad-well-pointedness-torsionless-epi}, but we write it out for completeness.
  From the triangle identity we know that $\varepsilon_{FX} \circ F\eta_X = \id_{FX}$.
  Consider the naturality square of $\varepsilon$ applied to $F\eta_X$:
  \begin{center}
    \begin{tikzcd}
      FGFX \ar{r}{FGF\eta_X} \ar{d}{\varepsilon_{FX}} & FGFGFX \ar{d}{\varepsilon_{FGFX}} \\
      FX \ar{r}{F\eta_X} & FGFX
    \end{tikzcd}
  \end{center}
  Combining this with the assumption $GF\eta_X = \eta_{GFX}$, we get that $F\eta_X \circ \varepsilon_{FX} = \varepsilon_{FGFX} \circ FGF\eta_X = \varepsilon_{FGFX} \circ F\eta_{GFX} = \id_{FGFX}$, where the last equality is given by the triangle identity.
  Thus $F\eta_X$ is an isomorphism, with inverse $\varepsilon_{FX}$.

  Now assume $F\eta_X$ is an isomorphism.
  By functoriality, $GF\eta_X$ is also an isomorphism with inverse $\mu_X$.
  Since $\eta_{GFX}$ also has $\mu_X$ as inverse, we get that $GF\eta_X = \eta_{GFX}$ by uniqueness of inverses.
\end{proof}

Let us apply the above discussion to the contravariant functor $(\argmhole)^\ast$.
This functor is right adjoint to itself, since $\Hom_R(A, B^\ast) \simeq \Hom_R(B,A^\ast)$.
To fit this into the above discussion, we take $F = (\argmhole)^\ast : \modcat{R} \to \modcat{R}^{\op}$
and $G = (\argmhole)^\ast : \modcat{R}^{\op} \to \modcat{R}$.
The unit and counit of this adjunction are both given by the evaluation map $\ev_A : A \to A^{\ast\ast}$.
Thus the associated monad is $T \coloneq (\argmhole)^{\ast\ast}$, with multiplication $\mu_A \coloneq (\ev_{A^\ast})^\ast$.
A module $X$ is $T$-reflexive if $\ev_X$ is an isomorphism, which is equivalent to $X$ being reflexive.
We say that a module $X$ is \notion{torsionless} if it is $T$-torsionless for this $T$,
i.e., if $\ev_X$ is a monomorphism.
This is equivalent to being torsionless in the sense of H. Bass~\cite{bass1960} and motivates the name.
We can specialize \Cref{prp:monad-well-pointedness-torsionless-epi} with the improvement from \Cref{prp:monad-adjunction-well-pointed-improvement} to this situation.

\begin{corollary}\label{cor:dbl-dual-reflexive-iss-dual-reflexive}
  Let $A$ be an $R$-module. The following are equivalent:
  \begin{enumerate}[noitemsep]
    \item $A^{\ast\ast}$ is reflexive
    \item $A^\ast$ is reflexive
    \item $\ev_A$ is torsionless-epic. \earlyqed
  \end{enumerate}
\end{corollary}

\begin{proposition}\label{prp:dbl-dual-tensor-factorization}
  Let $A$, $B$, and $C$ be finitely adequate modules. Then every map $f : A \otimes B \to C$ factors uniquely through $\ev_{A \otimes B} : A \otimes B \to A \dbldualtensor B$.
\end{proposition}

\begin{proof}
  By \Cref{prp:monad-factorization}, we see that $f$ factors uniquely through $\ev_{A \otimes B}$ if both $A \dbldualtensor B$ and $C$ are reflexive.
  The module $C$ is reflexive by \Cref{thm:can-dual-iso}, as it is finitely adequate.
  The double dual tensor product $A \dbldualtensor B$ is reflexive if $(A \otimes B)^\ast$ is reflexive by \Cref{cor:dbl-dual-reflexive-iss-dual-reflexive}.
  Now, $(A \otimes B)^\ast = \Hom_R(A \otimes B, R) \simeq \Hom_R(A, B^\ast)$, which is reflexive by \Cref{prp:fin-adeq-hom-closure} since $A$ and $B$ are finitely adequate.
\end{proof}

\begin{corollary}\label{cor:dbl-dual-tensor-adjunction}
  Let $M$ be a finitely adequate module.
  The functor $(\argmhole) \dbldualtensor M : \finadeq \to \finadeq$ is left adjoint to $\Hom_R(M,\argmhole) : \finadeq \to \finadeq$. \earlyqed
\end{corollary}
